\documentclass[11pt]{article}
\usepackage[a4paper,margin=1in]{geometry}
\usepackage{fontspec}
\usepackage[boldfont=false]{xeCJK}
\setCJKmainfont{FandolSong-Regular.otf}
\usepackage{amsmath}
\usepackage{amssymb}
\usepackage{array}
\usepackage{booktabs}
\usepackage{graphicx}
\usepackage{adjustbox}
\usepackage{enumitem}
\setlist[itemize]{topsep=2pt,partopsep=0pt,parsep=0pt,itemsep=2pt,leftmargin=1.4em}
\setlist[enumerate]{topsep=2pt,partopsep=0pt,parsep=0pt,itemsep=2pt,leftmargin=1.6em}
\usepackage{parskip}
\usepackage{fvextra}
\fvset{breaklines=true,breakanywhere=true,fontsize=\small}
\setlength{\parindent}{0pt}

\begin{document}

\begin{center}
  {\LARGE\bfseries Hydrocarbons}\\[0.35em]
  {\large A Level Chemistry}
\end{center}
\vspace{0.6em}

\section*{Alkanes}

\textbf{Alkanes} 烷烃 are saturated hydrocarbons (general formula $\text{C}_n\text{H}_{2n+2}$).

\subsection*{Making alkanes}

\begin{itemize}
\item \textbf{hydrogenation} 氢化: add hydrogen to an \textbf{alkene} 烯烃, using a $\text{Pt}$ or $\text{Ni}$ catalyst and heat.

\item \textbf{cracking} 裂化: break a long-chain alkane into shorter ones by heating with $\text{Al}_2\text{O}_3$.

\end{itemize}

\subsection*{Combustion}

In \textbf{combustion} 燃烧 an alkane burns in oxygen:

\begin{itemize}
\item \textbf{complete combustion} 完全燃烧 (plenty of oxygen) gives carbon dioxide and water.

\item \textbf{incomplete combustion} 不完全燃烧 (not enough oxygen) gives water plus toxic \textbf{carbon monoxide} 一氧化碳 ($\text{CO}$) and soot (carbon).

\end{itemize}

\subsection*{Free-radical substitution}

Alkanes react with chlorine or bromine by \textbf{free-radical substitution} 自由基取代, in \textbf{ultraviolet light} 紫外线. For ethane and chlorine the mechanism has three steps:

\begin{itemize}
\item \textbf{initiation} 引发 — UV light splits the halogen into two \textbf{free radicals} 自由基: $\;\text{Cl}_2 \rightarrow 2\,\text{Cl}\cdot$

\item \textbf{propagation} 增长 — radicals react and make new radicals: \[\text{Cl}\cdot + \text{C}_2\text{H}_6 \rightarrow \text{C}_2\text{H}_5\cdot + \text{HCl} \qquad \text{C}_2\text{H}_5\cdot + \text{Cl}_2 \rightarrow \text{C}_2\text{H}_5\text{Cl} + \text{Cl}\cdot\]

\item \textbf{termination} 终止 — two radicals join, ending the chain: $\;\text{Cl}\cdot + \text{C}_2\text{H}_5\cdot \rightarrow \text{C}_2\text{H}_5\text{Cl}$

\end{itemize}

\subsection*{Why cracking is useful, and why alkanes are unreactive}

Cracking turns heavy fractions of \textbf{crude oil} 原油 into more useful, lower-$M_r$ alkanes and alkenes. (A \textbf{fraction} 馏分 is a group of molecules with a similar boiling-point range.)

Alkanes are generally unreactive, especially towards polar reagents. This is because the C–H and C–C bonds are strong and have little \textbf{polarity} 极性, so there is no charge to attract an attacking species.

\subsection*{Environmental effects}

Burning alkanes in an internal combustion engine gives off carbon monoxide, oxides of nitrogen and unburnt hydrocarbons. A catalytic converter removes these by turning them into harmless gases.

\section*{Alkenes}

\textbf{Alkenes} have a C=C double bond (general formula $\text{C}_n\text{H}_{2n}$). The double bond is the functional group, so alkenes are reactive.

\subsection*{Making alkenes}

\begin{itemize}
\item \textbf{elimination} 消去 of $\text{HX}$ from a \textbf{halogenoalkane} 卤代烷, using $\text{NaOH}$ dissolved in ethanol, with heat.

\item \textbf{dehydration} 脱水 of an \textbf{alcohol} 醇, using a hot $\text{Al}_2\text{O}_3$ catalyst or concentrated sulfuric acid.

\item cracking of a longer-chain alkane.

\end{itemize}

\subsection*{Reactions of alkenes}

Most reactions are \textbf{electrophilic addition} 亲电加成 across the double bond:

\begin{center}
\adjustbox{max width=\linewidth}{%
\begin{tabular}{ll}
\toprule
\textbf{Reagent and conditions} & \textbf{Product} \\
\midrule
$\text{H}_2$, $\text{Pt/Ni}$ catalyst, heat & alkane \\
\textbf{steam} 水蒸气 ($\text{H}_2\text{O}$), $\text{H}_3\text{PO}_4$ catalyst & alcohol \\
a \textbf{hydrogen halide} 卤化氢 ($\text{HX}$), room temperature & halogenoalkane \\
a halogen $\text{X}_2$ & a di-substituted alkane \\
\bottomrule
\end{tabular}}
\end{center}

There are also two oxidation reactions with acidified $\text{KMnO}_4$:

\begin{itemize}
\item \textbf{cold, dilute} $\text{KMnO}_4$ adds two $\text{–OH}$ groups to give a \textbf{diol} 二醇.

\item \textbf{hot, concentrated} $\text{KMnO}_4$ breaks the C=C bond apart. The products (such as carboxylic acids or carbon dioxide) tell you where the double bond was.

\end{itemize}

\subsection*{Test for a C=C bond}

Shake the compound with orange \textbf{bromine water} 溴水. An alkene \textbf{decolourises} it (turns it colourless) by electrophilic addition. An alkane does not.

\subsection*{Addition polymerisation}

In \textbf{addition polymerisation} 加成聚合, many alkene molecules join into one long chain, with no other product. Ethene gives poly(ethene). The long-chain product is a \textbf{polymer} 聚合物.

\subsection*{The mechanism and Markovnikov's rule}

In electrophilic addition (for example bromine with ethene), the electron-rich C=C attracts the electrophile. This forms a positive intermediate called a \textbf{carbocation} 碳正离子, which the negative part then attacks.

Alkyl groups push electrons towards the positive carbon — this is the \textbf{inductive effect} 诱导效应. So a carbocation with more alkyl groups is more stable: tertiary is more stable than secondary, which is more stable than primary. When $\text{HBr}$ adds to propene, the more stable carbocation forms, so hydrogen adds to the carbon that already has more hydrogens. This pattern is \textbf{Markovnikov's rule} 马氏规则.


\end{document}
